\documentclass[11pt,letterpaper]{article}

% ── Geometry ──────────────────────────────────────────────────────────────────
\usepackage[margin=1in]{geometry}

% ── Pagination ────────────────────────────────────────────────────────────────
\usepackage[all]{nowidow}
\brokenpenalty=10000
\raggedbottom

% ── Fonts ─────────────────────────────────────────────────────────────────────
\usepackage[T1]{fontenc}
\usepackage{newpxtext}
\usepackage{newpxmath}

% ── Colors ────────────────────────────────────────────────────────────────────
\usepackage[dvipsnames]{xcolor}
\definecolor{SectionBlue}{HTML}{1B3A5C}
\definecolor{AccentGray}{HTML}{4A4A4A}
\definecolor{QuoteBorder}{HTML}{8B9DAF}
\definecolor{RiskRed}{HTML}{C0392B}
\definecolor{RiskAmber}{HTML}{E67E22}
\definecolor{RiskGreen}{HTML}{27AE60}
\definecolor{BoxBg}{HTML}{F7F8FA}

% ── Tables ────────────────────────────────────────────────────────────────────
\usepackage{booktabs}
\usepackage{tabularx}

% ── Boxes ─────────────────────────────────────────────────────────────────────
\usepackage[framemethod=tikz]{mdframed}

% ── Lists ─────────────────────────────────────────────────────────────────────
\usepackage{enumitem}
\usepackage{needspace}

% ── Product name (working title: change it here) ─────────────────────────────
\newcommand{\product}{Rooms}

% ── Hyperlinks ────────────────────────────────────────────────────────────────
\usepackage{hyperref}
\hypersetup{
  colorlinks=true,
  linkcolor=SectionBlue,
  urlcolor=SectionBlue,
  citecolor=SectionBlue,
  pdfauthor={Punt Labs},
  pdftitle={Rooms PR/FAQ: Lead a Team of Coding Agents Instead of Relaying Between Them},
  pdfsubject={Working Backwards: Hypothesis Stage},
}

% ── Section styling ───────────────────────────────────────────────────────────
\usepackage{titlesec}
\titleformat{\section}
  {\needspace{6\baselineskip}\Large\bfseries\color{SectionBlue}}{}{0pt}{}
\titleformat{\subsection}
  {\needspace{5\baselineskip}\large\bfseries\color{AccentGray}}{}{0pt}{}
\titleformat{\subsubsection}
  {\normalsize\bfseries\color{AccentGray}}{}{0pt}{}
\titlespacing*{\subsubsection}{0pt}{1em}{0.5em}[0pt]

% ── Document stage and version ───────────────────────────────────────────────
\newcommand{\prfaqstage}[1]{\def\theprfaqstage{#1}}
\prfaqstage{hypothesis}
\newcommand{\prfaqversion}[2]{\def\theprfaqmajor{#1}\def\theprfaqminor{#2}}
\prfaqversion{0}{1}

% ── Header/footer ────────────────────────────────────────────────────────────
\usepackage{fancyhdr}
\pagestyle{fancy}
\fancyhf{}
\fancyhead[L]{\footnotesize\textcolor{AccentGray}{Stage: \theprfaqstage{} \enspace\textbar\enspace v\theprfaqmajor.\theprfaqminor{} \enspace\textbar\enspace first draft}}
\fancyhead[R]{\footnotesize\textcolor{AccentGray}{\thepage}}
\renewcommand{\headrulewidth}{0pt}

% ══════════════════════════════════════════════════════════════════════════════
% Custom environments
% ══════════════════════════════════════════════════════════════════════════════

\newcommand{\prsection}[1]{%
  \vspace{1em}%
  {\large\bfseries\color{SectionBlue}#1}%
  \vspace{0.3em}\par%
}

\newmdenv[
  topline=false, bottomline=false, rightline=false,
  linewidth=3pt, linecolor=QuoteBorder, backgroundcolor=BoxBg,
  innerleftmargin=12pt, innerrightmargin=12pt,
  innertopmargin=8pt, innerbottommargin=8pt,
  skipabove=\baselineskip, skipbelow=\baselineskip,
]{customerquote}

\newmdenv[
  topline=false, bottomline=false, rightline=false,
  linewidth=3pt, linecolor=SectionBlue, backgroundcolor=BoxBg,
  innerleftmargin=12pt, innerrightmargin=12pt,
  innertopmargin=8pt, innerbottommargin=8pt,
  skipabove=\baselineskip, skipbelow=\baselineskip,
]{spokespersonquote}

\usepackage{changepage}

\newcounter{faqnum}
\newenvironment{faqpair}[1]{%
  \needspace{4\baselineskip}%
  \vspace{0.5em}%
  \refstepcounter{faqnum}%
  \noindent\textbf{\color{SectionBlue}Q\thefaqnum: #1}\par\nopagebreak[4]%
  \vspace{0.2em}%
  \begin{adjustwidth}{1em}{0pt}%
  \setlength{\parindent}{0pt}%
}{%
  \end{adjustwidth}%
  \vspace{0.5em}%
}
\newcommand{\faqref}[1]{\hyperref[#1]{FAQ~\ref*{#1}}}

\newcounter{featurenum}
\newcommand{\featureitem}[2]{%
  \refstepcounter{featurenum}%
  \item[\textbf{\color{SectionBlue}F\thefeaturenum.}]%
  \textbf{#1}: #2%
}
\newcommand{\featureref}[1]{\hyperref[#1]{Feature~\ref*{#1}}}

\setlength{\parskip}{0.5em}
\setlength{\parindent}{0pt}

% ══════════════════════════════════════════════════════════════════════════════
% Document
% ══════════════════════════════════════════════════════════════════════════════

\begin{document}

% ── Headline block ───────────────────────────────────────────────────────────
\begin{center}
  {\LARGE\bfseries\color{SectionBlue}
    \product{} Lets Developers Lead a Team of Coding Agents\\
    Instead of Relaying Messages Between Them} \\[0.8em]
  {\large\color{AccentGray}
    For the developer: one room per piece of work, spanning repos and machines.
    Agents bring you decisions, you answer by voice or a single key, and the
    state of the work is always one glance away.} \\[0.3em]
  {\large\color{AccentGray}
    For the agent: shared, typed state instead of chat scrollback, precise
    formats instead of paragraphs, and self-describing tools it can discover,
    use, and build on the spot.}
\end{center}

\vspace{0.5em}

% ══════════════════════════════════════════════════════════════════════════════
% PRESS RELEASE
% ══════════════════════════════════════════════════════════════════════════════

\section*{Press Release}

AUSTIN, TX, [LAUNCH DATE]. Punt Labs today released \product{}, a
keyboard-first, voice-forward workspace for developers who run several AI
coding agents at once. Terminal multiplexers built for agents, such as
herdr\cite{herdr}, put every agent's terminal on one screen. That is
convenient, but the developer is still the hub: reading each terminal to find
out what an agent needs, carrying answers back, and copying one agent's
findings into another's prompt. \product{} removes the developer from the
middle. Work lives in a \emph{room} that spans every repository and machine it
touches. Agents publish what they know to the room's shared state, co-design
with each other in precise formats, and bring the developer only the decisions
that need a human. The developer answers by voice or with one key, and goes
back to thinking.

\prsection{Problem}

A developer running three or four coding agents across several repositories
spends a surprising share of the day as a message router. Agent A needs to
know what agent B decided. Agent C is blocked on a question nobody has seen.
A change to a shared protocol touches three repos on two machines, and the
only view of it is three terminals. Free-form chat between agents makes this
worse, not better: prose is imprecise for agents and expensive for people to
read. Today's tools aggregate the terminals but leave the work of
coordination, the reading, relaying, and reconstructing of state, with the
human.

\prsection{Solution}

\product{} organizes work the way a lead runs a team, not the way a terminal
manager arranges windows. The developer thinks about four things, in order:
what needs me, where the work stands, who is doing what, and, only when
checking, the evidence. Hosts, terminals, branches, and processes still exist,
but they sit one keystroke below the work instead of being the work.

\begin{itemize}[leftmargin=1.5em, itemsep=0.3em]
  \item \textbf{Rooms.} A room groups the repositories, machines, agents, and
    state for one body of work. A protocol change touching three repos on two
    machines is one work item with one state, one CI picture, and one merge
    order.
  \item \textbf{Decisions, not interruptions.} A blocked agent posts a
    structured question with options and a recommendation. \product{} speaks
    it in one short sentence and, unless the developer has toggled it off,
    listens for the answer. ``Option two'' resumes the agent.
  \item \textbf{A shared board.} Agents publish plans, results, specifications,
    and decisions to the room's typed, signed, replayable state. New and
    restarted agents read the board instead of asking the developer to catch
    them up.
  \item \textbf{Agent channels.} Agents co-design with each other in schemas,
    tables, diffs, and formal notation, and escalate only real disagreements.
  \item \textbf{Self-describing tools.} The defaults are tools, and agents can
    make new ones on request. Every tool describes itself in typed terms, so
    every other agent in the room can discover and use it immediately.
  \item \textbf{Opinionated and fast.} One command to install, nothing to
    configure, every action one key, terminal first, with an optional GPU
    display for rich views.
\end{itemize}

\begin{spokespersonquote}
``I already get real leverage when my agents co-design with each other instead
of me shuttling between them. \product{} makes that the default. The question
I want the tool to answer is not `what are my terminals doing?' but `what does
this work need from me, and where does it stand?'''
\par\vspace{0.3em}\hfill--- Jim Freeman, Punt Labs
\end{spokespersonquote}

\begin{customerquote}
``I run four agents across a backend, a client, and a shared schema repo. Before,
I spent my mornings reading terminals to find out who was stuck. Now the room
tells me. Yesterday I answered three blocked agents by voice while reviewing a
design doc, and never opened a terminal.''
\par\vspace{0.3em}\hfill--- [ILLUSTRATIVE CUSTOMER], staff engineer
\end{customerquote}

\prsection{Getting Started}

One command installs \product{}. Running it in a workspace creates a room
from the repositories it finds, attaches to the agents already running,
reattaches any detached sessions, and enables voice if a microphone and Vox are
available. There is no wizard, no settings screen, and no account. Overrides
live in one plain-text file that nobody needs to open.

% ══════════════════════════════════════════════════════════════════════════════
% FAQ
% ══════════════════════════════════════════════════════════════════════════════

\newpage
\section*{Frequently Asked Questions}

\subsection*{External FAQs}

\begin{faqpair}{What is \product{}?}\label{faq:what}
A workspace for leading several AI coding agents across repositories and
machines. It keeps agent sessions alive, shows the state of the work, routes
agents' questions to you, and lets agents share state and tools with each
other. It works in any terminal, and adds a GPU window for rich views when one
is available.
\end{faqpair}

\begin{faqpair}{How is this different from herdr or tmux?}\label{faq:herdr}
tmux and herdr manage terminals, and herdr adds per-pane agent state on top
\cite{herdr,tmuxcontrol}. That is aggregation: you see everything, and you
still do the coordinating. \product{} matches the basics (persistent sessions,
detach and reattach, agent state, jump to the blocked agent, remote machines)
and then changes the unit of work from the terminal to the room:

\renewcommand{\arraystretch}{1.3}
\begin{tabularx}{\linewidth}{@{} l X X @{}}
\toprule
 & \textbf{Terminal multiplexer} & \textbf{\product{}} \\
\midrule
Home screen & Terminals & Decisions, board, team \\
Blocked agent & A flag; you read its pane & A structured question; you answer by voice or one key \\
Agent to agent & Send a prompt to another pane & Co-design in typed formats; escalate only disagreements \\
Shared context & Each agent's scrollback & One typed, replayable board per room \\
Multi-repo work & One terminal per repo & One work item across repos and hosts \\
\bottomrule
\end{tabularx}
\renewcommand{\arraystretch}{1.0}
\end{faqpair}

\begin{faqpair}{Do I need a GUI?}\label{faq:gui}
No. \product{} is terminal first and works over SSH. A GPU display adds rich
views (diffs, plots, node graphs, specifications) beside or inside the work,
but nothing requires it.
\end{faqpair}

\begin{faqpair}{How does voice work? Is the microphone always on?}\label{faq:voice}
When an agent asks a question, \product{} speaks it in one short sentence, then
opens the microphone for a few seconds and closes it when you stop talking.
Your answer is matched against the question's options, so recognition only
has to tell two or three choices apart. The microphone is closed while
\product{} speaks and at all other times. One key toggles listening off for a
call or a shared room; questions keep arriving on screen. Anything
irreversible is read back for confirmation. Every spoken answer is recorded as
text.
\end{faqpair}

\begin{faqpair}{What is a room?}\label{faq:room}
The container for one body of work: its repositories (on any of your
machines), its agents, its people, and its state. A repository can belong to
several rooms; each branch and session belongs to exactly one work item in one
room. Launching \product{} in a workspace creates a room from what it finds;
naming repositories by host and path creates one explicitly.
\end{faqpair}

\begin{faqpair}{Which agents does it work with?}\label{faq:agents}
Claude Code and Codex first, through hooks that declare an agent's state
(working, blocked, idle) rather than guessing it from terminal output. Any
agent that speaks MCP can read the board, raise decisions, and use the room's
tools; other terminal agents still get persistent sessions and
output-based state as a fallback.
\end{faqpair}

\begin{faqpair}{Agents can build tools. Can they break things?}\label{faq:tools-safety}
Tools come in two tiers. A tier-one tool is a composition: a query over the
room's state, a view built from a fixed set of elements, and actions that map
to existing operations. It contains no new code, so an agent can build one on
request and you keep or discard it with one key. A tier-two tool adds new code
and runs as a separate program, never inside the display, and installs only
after you approve its declared scope. Every tool, and every use of one, is a
signed, recorded fact.
\end{faqpair}

\begin{faqpair}{Does my code or data leave my machines?}\label{faq:local}
No. \product{} runs on machines you control. Rooms that span machines connect
them directly with mutual TLS; there is no hosted service and no account.
Agents send whatever they already send to their own model providers;
\product{} adds nothing to that.
\end{faqpair}

\begin{faqpair}{What does it cost?}\label{faq:cost}
\product{} is free and open source. [LICENSE TO BE DECIDED.]
\end{faqpair}

\subsection*{Internal FAQs}

\begin{faqpair}{Why now?}\label{faq:why-now}
Three things changed in 2025--2026. Developers routinely run several agents at
once, and agent multiplexers such as herdr gathered large followings quickly
\cite{herdr}, which shows the demand. Agents became good enough to co-design
with each other, which Punt Labs already sees with Biff. And the blackboard
architecture, shelved in the 1990s because its specialists were expensive
hand-built rule systems, became practical again now that capable general
specialists are cheap \cite{han2025bmas,salemi2025blackboard}.
\end{faqpair}

\begin{faqpair}{What is the user's mental model?}\label{faq:model}
A lead running a team. The explicit hierarchy is what the developer navigates;
the implicit hierarchy is the machinery, always one keystroke down and never
the starting point.

\renewcommand{\arraystretch}{1.3}
\begin{tabularx}{\linewidth}{@{} l X @{}}
\toprule
\textbf{Explicit} & Room $\rightarrow$ decisions, board, work items, team, tools \\
\textbf{Implicit} & Host $\rightarrow$ hub $\rightarrow$ repository $\rightarrow$ branch or worktree $\rightarrow$ session $\rightarrow$ PTY $\rightarrow$ terminal; plus commits, pull requests, CI runs \\
\textbf{The join} & The work item: an issue, its branches in one or more repos, their sessions, PRs, and CI \\
\bottomrule
\end{tabularx}
\renewcommand{\arraystretch}{1.0}

Agents are identities, not processes: a restarted session is the same team
member. Implicit failures surface in work terms (``two work items paused:
their agent's machine is offline''), never as raw infrastructure.
\end{faqpair}

\begin{faqpair}{Why shared state rather than chat?}\label{faq:shared-state}
Because coordination through shared, structured state scales where messaging
does not. The idea is old and well tested: Newell's 1962 blackboard metaphor
\cite{newell1962}, Hearsay-II's specialists posting hypotheses to a shared
structure \cite{erman1980hearsay}, Nii's surveys \cite{nii1986}, and Linda's
tuple spaces, which decouple writers and readers in space and time
\cite{gelernter1985linda}. Shared state lets a late-joining agent read where
things stand instead of replaying history; it tolerates missed updates, since
the next read corrects them; and its contents are machine-readable. Chat
remains, as a view of the same log.
\end{faqpair}

\begin{faqpair}{What data structure sits at the core?}\label{faq:facts}
One append-only log of typed facts, each a tuple of entity, attribute, value,
and transaction, with indexes for the current state. This is the model Datomic
popularized \cite{datomic}. Every transaction carries a sequence number and
the author's signed identity. The operations are Linda's plus time: assert
facts, query by pattern, take atomically (so two agents cannot claim the same
work), subscribe to a query's changes, and query the state as of any earlier
point. The board is the current state of the facts; the transcript is the log
read in order. Attributes are namespaced and typed, and each declares who may
write it. Terminal output, images, and large diffs are stored beside the log
and referenced from it, keeping the log small. The whole hub state is a
relation, which suits the Z specifications Punt Labs already uses to verify
concurrency.
\end{faqpair}

\begin{faqpair}{How does a room span repositories and machines?}\label{faq:mesh}
Each room has a home hub that sequences its transactions; other machines that
hold a member repository or run a member agent replicate the log and join
implicitly. Write rate is human and agent scale, hundreds of facts per second
at most, so one sequencer per room is not a bottleneck. A machine cut off from
the home hub keeps working and submits queued writes on reconnection, where
the write rules are applied. That reconnection path is the part to
model-check before it is built. Machines connect with mutual TLS against a
pluggable trust anchor.
\end{faqpair}

\begin{faqpair}{How are tools defined and discovered?}\label{faq:tools}
A tool is itself a set of facts: its name, purpose, what it reads, what it
writes, its typed operations, its view, its key binding, its voice phrase, and
its scope. Discovery is a query; an agent can also subscribe to new tools. The
room's MCP endpoint publishes tools as MCP tools and notifies clients when the
list changes \cite{mcplistchanged}, so agents can use a new tool without
restarting. The defaults (decisions, board, work items, team) are built on the
same interface as any other tool. This is the Smalltalk principle: if the
built-ins cannot be written with the public interface, the interface is
incomplete.
\end{faqpair}

\begin{faqpair}{Why merge Lux, Vox, and Biff?}\label{faq:merge}
They are three projections of one goal: communication between people and
agents that is more precise and less tiring than free-form text. Lux renders
state visually, Vox speaks and listens, and Biff carries messages, presence,
and identity across machines. \product{} defines one artifact protocol and one
state model; Lux becomes the visual client, Vox the voice client, and Biff the
transport. The three products keep working while the shared protocol takes
shape; the code merges later, if at all.
\end{faqpair}

\begin{faqpair}{Why a Rust or Go core, and why terminal first?}\label{faq:stack}
The hub is a long-running concurrent daemon, the kind of program where a
compiled language with strong concurrency guarantees pays off most; much of
Lux's recent design work fixed races and render-loop cost in Python. A single
static binary is also simpler to install. Terminal first because the
strongest competitor runs anywhere a terminal does, including SSH, and a
GPU-only product would lose on that alone. Lux's first external integration
was lost partly to protocol drift (client at 0.1, Lux at 0.33), which is why
\product{} treats a small, versioned protocol with a published schema and a
conformance suite as the product.
\end{faqpair}

\begin{faqpair}{What do we reuse?}\label{faq:reuse}
Lux's element vocabulary, self-validation, introspection, applet model, and
mutual-TLS transport design; its formal models and their method; Vox's
synthesis, notification levels, and voice-input spike; Biff's messaging,
presence, and identity; ethos identities and missions; Linear as the source of
work items. The conformance suite is built from Lux's end-to-end scenarios.
\end{faqpair}

\begin{faqpair}{How do we measure success?}\label{faq:metrics}
By what the developer no longer does, compared with a terminal multiplexer on
the same workload:

\begin{itemize}[leftmargin=1.5em, itemsep=0.2em]
  \item relay hops per hour (the developer carrying information between agents);
  \item interruptions per hour, and the share that needed a human;
  \item time from an agent blocking to its question being answered;
  \item share of questions answered without opening a terminal;
  \item time from install to the first answered question (target: one command, no decisions).
\end{itemize}

Targets are set before the first external trial, not after.
\end{faqpair}

\begin{faqpair}{What happens to generative UI?}\label{faq:genui}
It stays constrained, which is what makes it fast and trustworthy. Agents
compose views from a fixed element vocabulary inside a stable frame; they do
not generate code for the renderer. Research shows people prefer generated
interfaces to text but flags latency and inconsistency
\cite{leviathan2026generativeui}; a fixed frame with typed slots answers
consistency, and preparing the next screens from an enumerable set of
responses (a decision's options) answers latency.
\end{faqpair}

\begin{faqpair}{What is the plan?}\label{faq:plan}
\begin{enumerate}[leftmargin=1.5em, itemsep=0.2em]
  \item Protocol schema and conformance suite; the fact log and room sequencer, model-checked.
  \item Session host and terminal client at parity with the basics: persistent sessions, agent state through hooks, jump to blocked.
  \item Decision requests with voice output and listen-after-questions. This is the first thing users should feel.
  \item Agent channels through Biff, and the board.
  \item Rooms across machines.
  \item Tools: tier one by agents on request, then tier two.
  \item The GPU client for rich views.
\end{enumerate}
\end{faqpair}

\begin{faqpair}{What are the open questions?}\label{faq:open}
\begin{itemize}[leftmargin=1.5em, itemsep=0.2em]
  \item The product name, and whether it ships under the Lux name.
  \item Rust or Go for the core.
  \item Whether a room's manifest lives on its home hub only, or also as a file in a repository.
  \item Whether the board lives with Biff (transport and identity) or with the hub.
  \item How competing proposals with confidence, Hearsay-II's strongest idea, appear on the board.
  \item The license.
\end{itemize}
\end{faqpair}

% ══════════════════════════════════════════════════════════════════════════════
% RISK ASSESSMENT
% ══════════════════════════════════════════════════════════════════════════════

\newpage
\section*{Risk Assessment}

\renewcommand{\arraystretch}{1.6}
\begin{tabularx}{\textwidth}{@{} l l X @{}}
\toprule
\textbf{\color{SectionBlue}Risk} & \textbf{\color{SectionBlue}Rating} & \textbf{\color{SectionBlue}Assessment} \\
\midrule
Value & \textcolor{RiskAmber}{Medium} &
  The core claim, that structured asks, a shared board, and agent channels
  reduce the developer's routing and interruptions, rests on one developer's
  experience with Biff. No external user has tried it. The metrics in
  \faqref{faq:metrics} are the test. \\
Usability & \textcolor{RiskAmber}{Medium} &
  Voice must be reliable enough to trust: false triggers, missed answers,
  and noisy rooms all erode it. Constrained answers and the one-key toggle
  mitigate this, but only an external trial measures it. A keyboard-first
  terminal interface has a learning curve some users will not climb. \\
Feasibility & \textcolor{RiskRed}{High} &
  The scope is large: a new core in a new language, a terminal emulator good
  enough to compete, a replicated fact log with offline reconnection, voice
  input, and three products' integration. The stable protocol and conformance
  suite reduce the risk of moving off Python but not the size of the work. \\
Viability & \textcolor{RiskAmber}{Medium} &
  Free and open source, with value as the centre of the Punt Labs tool suite.
  herdr's head start and momentum mean \product{} wins only by leading with
  what herdr structurally cannot do, not by parity. \\
\bottomrule
\end{tabularx}
\renewcommand{\arraystretch}{1.0}

% ══════════════════════════════════════════════════════════════════════════════
% FEATURE APPENDIX
% ══════════════════════════════════════════════════════════════════════════════

\section*{Feature Appendix}

\subsection*{Must Do}

\begin{enumerate}[leftmargin=2.5em, itemsep=0.3em]
  \featureitem{Rooms}{a container for repositories on any machine, agents, people, state, and tools, created from the launch workspace by default}\label{feat:rooms}
  \featureitem{Persistent sessions}{agent sessions held by a session host, surviving client exit, detach, and reattach}
  \featureitem{Declared agent state}{working, blocked, and idle reported through hooks, with output matching as a fallback}
  \featureitem{Decision requests}{structured questions with options and a recommendation, answered with one key}\label{feat:decisions}
  \featureitem{Voice}{one-sentence spoken questions, listen-after-questions, constrained answers, read-back for irreversible actions, and a one-key toggle}\label{feat:voice}
  \featureitem{The board}{the room's typed, signed, replayable shared state}
  \featureitem{Cross-repo work items}{one work item spanning branches in several repositories, with one state, combined CI, and a merge order}
  \featureitem{Keyboard layer}{every action one key, with automatic tiling and no required mouse}
  \featureitem{Zero configuration}{one command from install to the first answered question}
  \featureitem{Terminal client}{the full product in any terminal, including over SSH}
\end{enumerate}

\subsection*{Should Do}

\begin{enumerate}[leftmargin=2.5em, itemsep=0.3em, resume*]
  \featureitem{Agent channels}{agent-to-agent co-design in typed formats through Biff, escalating only disagreements}
  \featureitem{Agent-built tools}{tier-one tools composed on request, self-describing and discoverable by every agent}\label{feat:tools}
  \featureitem{Rooms across machines}{replication, implicit host membership, and offline reconnection}
  \featureitem{GPU client}{rich views beside and inside the work: diffs, plots, node graphs, specifications}
  \featureitem{Facts inspector and replay}{the room's state at any earlier point, and every write's author}
\end{enumerate}

\subsection*{Won't Do}

\begin{enumerate}[leftmargin=2.5em, itemsep=0.3em, resume*]
  \featureitem{Hosted service or accounts}{\product{} runs on machines the user controls}
  \featureitem{Arbitrary code execution in the core}{new code runs only as scoped, approved tier-two tools outside the display}
  \featureitem{A setup wizard or settings screen}{defaults are the product}
  \featureitem{Hosts as the primary navigation}{machines appear only when something fails}
  \featureitem{Chat-first coordination}{chat is a view of the log, not the coordination mechanism}
  \featureitem{An extension marketplace}{tools are built per room, by its people and agents}
\end{enumerate}

% ══════════════════════════════════════════════════════════════════════════════
% REFERENCES
% ══════════════════════════════════════════════════════════════════════════════

\begin{thebibliography}{99}

\bibitem{herdr}
herdr: a terminal multiplexer for AI coding agents.
\url{https://github.com/herdrdev/herdr}. Accessed November 2026.

\bibitem{tmuxcontrol}
tmux wiki, ``Control Mode.''
\url{https://github.com/tmux/tmux/wiki/Control-Mode}.

\bibitem{newell1962}
A.~Newell, ``Some Problems of Basic Organization in Problem-Solving Programs,''
in \emph{Self-Organizing Systems}, 1962.

\bibitem{erman1980hearsay}
L.~D. Erman, F.~Hayes-Roth, V.~R. Lesser, and D.~R. Reddy, ``The Hearsay-II
Speech-Understanding System: Integrating Knowledge to Resolve Uncertainty,''
\emph{ACM Computing Surveys} 12(2), 1980.

\bibitem{nii1986}
H.~P. Nii, ``Blackboard Systems,'' parts 1 and 2, \emph{AI Magazine} 7(2--3),
1986.

\bibitem{gelernter1985linda}
D.~Gelernter, ``Generative Communication in Linda,'' \emph{ACM Transactions on
Programming Languages and Systems} 7(1), 1985.

\bibitem{datomic}
Datomic, ``Datomic Information Model.'' \url{https://docs.datomic.com}.

\bibitem{han2025bmas}
B.~Han and S.~Zhang, ``Exploring Advanced LLM Multi-Agent Systems Based on
Blackboard Architecture,'' arXiv:2507.01701, 2025.

\bibitem{salemi2025blackboard}
A.~Salemi et al., ``LLM-Based Multi-Agent Blackboard System for Information
Discovery in Data Science,'' arXiv:2510.01285, 2025.

\bibitem{mcplistchanged}
Model Context Protocol specification, ``Tools: list changed notification.''
\url{https://modelcontextprotocol.io}.

\bibitem{leviathan2026generativeui}
Y.~Leviathan et al., ``Generative UI: LLMs are Effective UI Generators,''
arXiv:2604.09577, 2026.

\end{thebibliography}

\end{document}
